\section{Discussion}\label{sec:discussion}

\subsection{Are machine learning improvements statistically real?}

Our results give a nuanced answer to the question in the title. With spatially or temporally dependent data, the conventional pipeline (random $K$-fold cross-validation followed by a paired $t$-test, an i.i.d. bootstrap, or even the corrected resampled $t$-test) produces false discoveries at rates between $15\%$ and $40\%$ under moderate to strong dependence in our simulations, instead of the nominal $5\%$. Moreover, random cross-validation targets an interpolation estimand that can understate the gain at new places and times by half, and the case studies show that the ranking of the same learners can change with the prediction task: on hourly weather data, XGBoost appears best under random folds, the learners are indistinguishable when filling two-week gaps, and the random forest is significantly better when forecasting later weeks. Many reported ``significant'' improvements on such data are therefore not statistically supported, and some are answers to a different question from the one of interest.

Switching to blocked cross-validation fixes the estimand but not, by itself, the inference: applied to blocked folds, procedures that resample observations or that model only temporal dependence are \emph{more} anti-conservative than under random folds ($60$--$80\%$ false positives under strong dependence). Cluster procedures (block bootstrap, random-intercept models) remain anti-conservative whenever the dependence range is not small compared with the blocks, and they ignore the correlation of training-induced fold effects. The blocked paired $t$-test, whose behaviour is explained by Proposition~\ref{prop:pairedt}, is the best conventional procedure but still rejects $11$--$18\%$ of true nulls under moderate and strong dependence.

\subsection{What DA-P adds, and what it does not}

\paragraph{Relative to the Hering--Genton test.} The closest existing procedure is the spatial prediction comparison test of Hering and Genton~\cite{hering2011}, which we extended to space--time. When the training-induced variability is small (Experiment~1 and its variants), both procedures are valid, and Hering--Genton is somewhat more powerful because DA-P spends degrees of freedom on a training term that is nearly zero. When the models are refitted learners with substantial fold-to-fold variability (Experiment~2), Hering--Genton rejects $17$--$23\%$ of true nulls under moderate and strong dependence, against $10$--$11\%$ for DA-P; on real data, it declares one seed of XGBoost significantly better than another in $20\%$ of the pedestrian pairs (DA-P: $3\%$), and reports a significant forecasting difference between XGBoost and logistic regression that DA-P attributes to fold-to-fold variability. The training term, whose estimation and limits are characterised by Proposition~\ref{prop:bias}, is thus what distinguishes the comparison of \emph{learning procedures} from the comparison of fixed predictions.

\paragraph{Relative to the kernel variant.} Our initial kernel version, DA-K, over-rejected under moderate and strong dependence ($10$--$15\%$); the pilot study traced this to replicates in which the training variance was truncated at zero, where the downward bias of the fold-centred kernel estimate was no longer compensated. The parametric variance of DA-P does not suffer from kernel down-weighting, and it remained valid under anisotropy, a smoother covariance, heteroscedasticity, non-stationary ranges and irregular sampling, misspecification costing power rather than validity.

\paragraph{What remains.} DA-P is not exactly calibrated in the most difficult setting we studied: learners compared under moderate or strong dependence, where it rejects $10$--$11\%$ of true nulls ($8.7\%$ with the three-way rule). Proposition~\ref{prop:bias} identifies the reason: whichever covariate field happens to be more informative over the whole sample favours one model in every fold, which acts as a fold effect common to all folds, i.e.\ a training correlation $\rho$ larger than assumed, and such a component cannot be estimated from a single cross-validation run~\cite{bengio2004}. The same mechanism explains the over-rejection of all procedures when fold effects are highly correlated in Section~\ref{sec:sim1rho}. The range ratio, which we initially intended to use as a guard, does not identify these cases (Section~\ref{sec:sim1rule}); it is a useful indicator of how much information a comparison can carry, not of its validity.

\subsection{Recommendations}

From these results we derive the following recommendations for comparing two models on spatiotemporal data.
\begin{enumerate}
\item \emph{Define the estimand first.} Choose the validation design to match the intended use of the model: random cross-validation answers an interpolation question at the scale of single observations, blocked cross-validation an interpolation question at the scale of the blocks, and forward chaining a forecasting question.
\item \emph{Never rely on procedures that treat observations as independent} (i.i.d. bootstrap, $z$- or $t$-tests on pooled observations, McNemar's test) or on the paired $t$-test with random folds. Under dependence, the corrected resampled $t$-test is not a remedy either.
\item \emph{Use blocked cross-validation with a test that accounts for both sources of dependence}, such as DA-P, and report its result under $\rho_0=1/K$ and $\rho_0=0.5$; claim an improvement only when it holds under both (the three-way rule), and report the other cases as indeterminate.
\item \emph{Inspect the diagnostics.} Report the estimated ranges, the effective sample size and the range ratio. A small effective sample (a few dozen, as on the Las Rosas field) means that only large differences can be resolved.
\item \emph{Correct for multiplicity} when several models are compared, and treat comparisons on a single data set as exploratory.
\item \emph{Report intervals, not only $p$-values}, on the scale of the loss, as advocated for large data sets in this journal~\cite{caponero2026}: an improvement can be statistically real and practically negligible, or practically important and statistically unresolved.
\end{enumerate}

\subsection{Limitations and future work}

Our simulation designs are stylised: Gaussian fields on a $12\times12$ or $10\times10$ grid (with one irregular variant), a single learner family in Experiment~2, and $K=5$ folds throughout. The working model~\eqref{eq:model} assumes additive fold effects and a stationary loss differential, which may fail when one model is better in some regions and worse in others; a version with a spatially varying $\Delta$ is a natural extension. The correlation $\rho$ of training effects is assumed rather than estimated; repeated blocked cross-validation with different block allocations, or nested cross-validation in the spirit of Bates et al.~\cite{bates2024}, could provide information on it and is, in our view, the most promising route to a fully calibrated test. The exponential variogram does not represent periodic dependence such as daily and weekly cycles, which matter for the pedestrian data. We considered additive losses only; non-additive metrics such as the AUC or the F1-score would require linearisation. Finally, the choice of the parametric variance and of the three-way rule was informed by pilot studies on separate seeds and, for the rule, by the failure of a range-based alternative on the main replicates; both deserve confirmation on further designs and data sets.
